\documentclass[11pt,a4paper]{article}
\usepackage[T1]{fontenc}
\usepackage[utf8]{inputenc}
\usepackage{amsmath,amssymb,amsthm}
\usepackage[margin=1in]{geometry}
\usepackage[colorlinks=true,linkcolor=blue,urlcolor=blue]{hyperref}
\theoremstyle{plain}
\newtheorem{theorem}{Theorem}
\newtheorem{lemma}[theorem]{Lemma}
\newtheorem{proposition}[theorem]{Proposition}
\newtheorem{corollary}[theorem]{Corollary}
\theoremstyle{definition}
\newtheorem{remark}[theorem]{Remark}
% A quoted conjecture can be a single hundred-term formula whose terms are thirty-digit
% numbers. TeX cannot hyphenate those, so without a little stretch every such quote runs off
% the page; the linked-a-file papers are all of that shape.
\setlength{\emergencystretch}{3em}

\title{The empirical recurrence for OEIS A234121, proved by transfer matrix}
\author{Adrian Perez Fontelles\\ \small Independent researcher}
\date{14 September 2026}
\begin{document}
\maketitle

\begin{abstract}
OEIS A234121 counts $(n+1)\times8$ arrays over $\{0,\dots,2\}$ subject to a condition imposed on
every $2\times2$ subblock, and carries an empirical recurrence of order $62$ contributed
by R.~H.~Hardin. It is true, and it is not an empirical matter at all. The condition is local
to each subblock, so the rows of an admissible array are the vertices of a finite digraph
on $S=6561$ vertices whose edges are the admissible consecutive pairs, and the count is a
number of walks: $a(n)=\mathbf{1}^{\!\top}M^{n}\mathbf{1}$. Writing $q$
for the characteristic polynomial of the conjectured recurrence, the recurrence holds for all
$n$ exactly when $\mathbf{1}^{\!\top}M^{j}q(M)\mathbf{1}=0$ for every $j\ge0$, and the
Cayley--Hamilton theorem bounds that to $j<S$. It is therefore a finite computation in exact
integer arithmetic, and it comes out zero.
\end{abstract}

\noindent\small 2020 Mathematics Subject Classification. 05A15, 05B45, 15B36.\normalsize

\section{The sequence and the conjecture}

OEIS A234121 is ``Number of (n+1)X(7+1) 0..2 arrays with every 2X2 subblock having the absolute values of all six edge and diagonal differences no larger than 1.''. It has offset $1$ and begins
\[
313537, 93113761, 26937444801, 7888454356625, 2315408571668225, 680845368964238145, 200376051988817866625, 58999673293833047076529,\ \dots
\]
The entry carries the following comment:
\begin{quote}\small
Empirical recurrence of order 62 (see link above)

Empirical: a(n) = 553*a(n-1) - 85548*a(n-2) + 1053884*a(n-3) + 596445686*a(n-4) - 20859621734*a(n-5) - 1838422509868*a(n-6) + 64740000330460*a(n-7) + 3452743975554537*a(n-8) - 87552711729579809*a(n-9) - 4172638342894466264*a(n-10) + 56233787710584949752*a(n-11) + 3150367371827445840196*a(n-12) - 13153320428690338732756*a(n-13) - 1442082901267441502521144*a(n-14) - 2960086642834181354168168*a(n-15) + 399183922871796742043588993*a(n-16) + 2563188501502610801542563191*a(n-17) - 66683442153378596122249439164*a(n-18) - 680375931772265186762950694516*a(n-19) + 6551888243909391729670099382326*a(n-20) + 98593727769464441508189852071034*a(n-21) - 338885566496324779194056711137500*a(n-22) - 8868137854363101424865827734715444*a(n-23) + 3152027495423783009884873488115959*a(n-24) + 528420352884761244733585078437057761*a(n-25) + 749813824706492157035990450334084416*a(n-26) - 21808774479101761396082555083053176000*a(n-27) - 55111703179168534460794572766032372752*a(n-28) + 646413544468142817561155458472447958112*a(n-29) + 2079682315656590071205610939878981264832*a(n-30) - 14228581355755613156971897656846109159616*a(n-31) - 50659200194969765578735390637355420783552*a(n-32) + 240146746400130039879934061430340928122368*a(n-33) + 850434674401514888727929154600339200059648*a(n-34) - 3188771825739453288515854718750150700581888*a(n-35) - 10036132272374181009512541990114342336148480*a(n-36) + 33629053735096807934237436372838776680054784*a(n-37) + 82571632452813033382575249848054889880473600*a(n-38) - 277623539113624015479637953996439996648734720*a(n-39) - 454216070172926849783614826844771991117332480*a(n-40) + 1729882740224894035589578264900951750149144576*a(n-41) + 1487413757723846440632660919774257270734979072*a(n-42) - 7739211556416346357151583926107969659353432064*a(n-43) - 1685718977153011086985694928645043283357073408*a(n-44) + 23475240168546693003507178743742855386771750912*a(n-45) - 6566397316402662625516677227364614714022690816*a(n-46) - 45238814137679251304779247785674538235040104448*a(n-47) + 30513008673693369348412413585915155060782792704*a(n-48) + 50828684715320496781113128946769515891325403136*a(n-49) - 53399210695796900387918972951032689034040180736*a(n-50) - 28895302367322499155375590074635642053603622912*a(n-51) + 47148895660849721538302624305588321919441043456*a(n-52) + 5778001374266269571245906668978765749878259712*a(n-53) - 22569485161171394895658677095383069386099130368*a(n-54) + 938008510912498923951720935594532099676176384*a(n-55) + 6150687945509169934795794363992477802172514304*a(n-56) - 471670794591749688007589944447477057780187136*a(n-57) - 963502964636557544404606443605590680847515648*a(n-58) + 29628618729595279540156195667851500085837824*a(n-59) + 74911313861590721218880701015473888228802560*a(n-60) + 3886298604222860830967627071661951208652800*a(n-61) - 1092304090113260063682602502675983499264000*a(n-62)
\end{quote}
As of the ``Last modified'' line on the live entry (Jun 02 2025 09:12:15, revision 7) the statement is
still recorded as empirical, and nothing on the entry records it as settled either way.

Written out, the claim is
\begin{equation}\label{eq:conj}
a(n)\;=\;553\,a(n-1) - 85548\,a(n-2) + 1053884\,a(n-3) + 596445686\,a(n-4) - 20859621734\,a(n-5) - 1838422509868\,a(n-6) + 64740000330460\,a(n-7) + 3452743975554537\,a(n-8) - 87552711729579809\,a(n-9) - 4172638342894466264\,a(n-10) + 56233787710584949752\,a(n-11) + 3150367371827445840196\,a(n-12) - 13153320428690338732756\,a(n-13) - 1442082901267441502521144\,a(n-14) - 2960086642834181354168168\,a(n-15) + 399183922871796742043588993\,a(n-16) + 2563188501502610801542563191\,a(n-17) - 66683442153378596122249439164\,a(n-18) - 680375931772265186762950694516\,a(n-19) + 6551888243909391729670099382326\,a(n-20) + 98593727769464441508189852071034\,a(n-21) - 338885566496324779194056711137500\,a(n-22) - 8868137854363101424865827734715444\,a(n-23) + 3152027495423783009884873488115959\,a(n-24) + 528420352884761244733585078437057761\,a(n-25) + 749813824706492157035990450334084416\,a(n-26) - 21808774479101761396082555083053176000\,a(n-27) - 55111703179168534460794572766032372752\,a(n-28) + 646413544468142817561155458472447958112\,a(n-29) + 2079682315656590071205610939878981264832\,a(n-30) - 14228581355755613156971897656846109159616\,a(n-31) - 50659200194969765578735390637355420783552\,a(n-32) + 240146746400130039879934061430340928122368\,a(n-33) + 850434674401514888727929154600339200059648\,a(n-34) - 3188771825739453288515854718750150700581888\,a(n-35) - 10036132272374181009512541990114342336148480\,a(n-36) + 33629053735096807934237436372838776680054784\,a(n-37) + 82571632452813033382575249848054889880473600\,a(n-38) - 277623539113624015479637953996439996648734720\,a(n-39) - 454216070172926849783614826844771991117332480\,a(n-40) + 1729882740224894035589578264900951750149144576\,a(n-41) + 1487413757723846440632660919774257270734979072\,a(n-42) - 7739211556416346357151583926107969659353432064\,a(n-43) - 1685718977153011086985694928645043283357073408\,a(n-44) + 23475240168546693003507178743742855386771750912\,a(n-45) - 6566397316402662625516677227364614714022690816\,a(n-46) - 45238814137679251304779247785674538235040104448\,a(n-47) + 30513008673693369348412413585915155060782792704\,a(n-48) + 50828684715320496781113128946769515891325403136\,a(n-49) - 53399210695796900387918972951032689034040180736\,a(n-50) - 28895302367322499155375590074635642053603622912\,a(n-51) + 47148895660849721538302624305588321919441043456\,a(n-52) + 5778001374266269571245906668978765749878259712\,a(n-53) - 22569485161171394895658677095383069386099130368\,a(n-54) + 938008510912498923951720935594532099676176384\,a(n-55) + 6150687945509169934795794363992477802172514304\,a(n-56) - 471670794591749688007589944447477057780187136\,a(n-57) - 963502964636557544404606443605590680847515648\,a(n-58) + 29628618729595279540156195667851500085837824\,a(n-59) + 74911313861590721218880701015473888228802560\,a(n-60) + 3886298604222860830967627071661951208652800\,a(n-61) - 1092304090113260063682602502675983499264000\,a(n-62)
\end{equation}
for all $n>62$.

\section{The condition is local, so the rows form a digraph}

Write an array of the shape in question as its list of rows
$r^{(1)},r^{(2)},\dots$, each an element of
$V=\{0,1,\dots,2\}^{8}$, so $|V|=S=6561$. A $2\times2$ subblock of the array
occupies two consecutive rows and two consecutive columns; if the two rows are $r$
and $s$ (in that order) and the two columns are numbered $j,j+1$,
the subblock is
\[
\begin{pmatrix}\alpha&\beta\\\gamma&\delta\end{pmatrix}=\begin{pmatrix}r_j&r_{j+1}\\s_j&s_{j+1}\end{pmatrix}.
\]
The entry's requirement on that subblock is
\begin{equation}\label{eq:loc}
\max\bigl(|\beta-\alpha|,|\delta-\beta|,|\gamma-\delta|,|\alpha-\gamma|,|\delta-\alpha|,|\gamma-\beta|\bigr)\;\le\;1,
\end{equation}
a condition on the four entries alone.

For $r,s\in V$ declare $r\to s$ an edge of a digraph $G$ when, for every
$1\le j\le 8-1$,
\begin{itemize}
\item the four entries above satisfy \eqref{eq:loc};

\end{itemize}
Let $M\in\{0,1\}^{S\times S}$ be the adjacency matrix of $G$ and $\mathbf{1}$ the
all-ones vector.

\begin{lemma}\label{lem:walk}
The arrays counted by the entry with $L$ rows are exactly the walks
$r^{(1)}\to r^{(2)}\to\cdots\to r^{(L)}$ of length $L-1$ in $G$. Consequently
\[
a(n)\;=\;\mathbf{1}^{\!\top}M^{n}\mathbf{1} .
\]
\end{lemma}

\begin{proof}
An array with $L$ rows and $8$ columns is precisely a sequence
$r^{(1)},\dots,r^{(L)}$ of elements of $V$; conversely any such sequence is an array.
Every $2\times2$ subblock lies in two consecutive rows $r^{(t)},r^{(t+1)}$ and two
consecutive columns, so the array satisfies the entry's condition on all of its subblocks
if and only if each consecutive pair $\bigl(r^{(t)},r^{(t+1)}\bigr)$ satisfies it for
every column pair --- which is exactly the edge relation of $G$. The number of walks of
length $L-1$, summed over all starting and ending vertices, is
$\mathbf{1}^{\!\top}M^{L-1}\mathbf{1}$, since $(M^{k})_{r,s}$ counts the walks of
length $k$ from $r$ to $s$. The entry's index $n$ corresponds to $L=n+1$ rows.
\end{proof}

\section{The criterion}

Let $q(t)=t^{62}-\sum_i c_i\,t^{62-i}$ be the characteristic polynomial of
\eqref{eq:conj}, with the $c_i$ read off that equation.

\begin{lemma}\label{lem:crit}
For every $n$ with $n+1-1\ge62$,
\[
a(n)-\sum_i c_i\,a(n-i)\;=\;\mathbf{1}^{\!\top}M^{\,n+1-1-62}\,q(M)\,\mathbf{1} .
\]
Hence \eqref{eq:conj} holds from that point on if and only if
$\mathbf{1}^{\!\top}M^{j}q(M)\mathbf{1}=0$ for every $j\ge0$; and it is enough to check
$j=0,1,\dots,S-1$.
\end{lemma}

\begin{proof}
By Lemma~\ref{lem:walk}, $a(n-i)=\mathbf{1}^{\!\top}M^{\,n+1-1-i}\mathbf{1}$, so
\[
a(n)-\sum_i c_i a(n-i)
=\mathbf{1}^{\!\top}\Bigl(M^{m}-\sum_i c_i M^{m-i}\Bigr)\mathbf{1}
=\mathbf{1}^{\!\top}M^{\,m-62}\,q(M)\,\mathbf{1}
\]
with $m=n+1-1$, which is the stated identity; the equivalence follows by letting
$m-62=j$ run over $\ge0$. For the last clause put $w=q(M)\mathbf{1}$. By the
Cayley--Hamilton theorem $M^{S}$ is an integer combination of $I,M,\dots,M^{S-1}$, so
$\mathbf{1}^{\!\top}M^{j}w$ for $j\ge S$ is the corresponding combination of the earlier
values; if those all vanish, so do all the rest.
\end{proof}

\section{The computation}

\begin{theorem}
The recurrence \eqref{eq:conj} holds for every $n>62$, which is the statement of
Section~1.
\end{theorem}

\begin{proof}
The digraph was constructed from \eqref{eq:loc} by a depth-first scan across the $8$
columns: the edge condition is a conjunction of one constraint per consecutive column
pair, so the successors of a row $r$ are enumerated column by column without testing all
$S^2$ pairs. The vector $w=q(M)\mathbf{1}\in\mathbb{Z}^{S}$ was then formed by $62$
matrix--vector products in exact integer arithmetic, and $\mathbf{1}^{\!\top}M^{j}w$ was
evaluated for $j=0,1,\dots$, again exactly. Every one of those integers vanishes from the
index corresponding to $n=62+1$ onwards, and the run continued until $w$ itself was the
zero vector, which settles every larger $j$ at once. By Lemma~\ref{lem:crit} the recurrence
holds for every $n>62$.
\end{proof}

\begin{remark}
The argument gives more than the recurrence: it shows the sequence is $C$-finite with
characteristic polynomial dividing that of $M$, so it has a rational generating function
whose denominator divides $\det(I-xM)$. The conjectured recurrence is one consequence; the
minimal one is a divisor of $q$.
\end{remark}

\section{Verification}

Three checks, each independent of the algebra above.

First, the digraph was built directly from the entry's stated condition and the walk counts
$\mathbf{1}^{\!\top}M^{L-1}\mathbf{1}$ compared against the entry's DATA at
every one of the $10$ published indices; the values agree exactly. That is what ties
the digraph to the sequence: a model that reproduced only some of the published terms would
be the wrong model, and would have been rejected. In particular it is what rules out a
misreading of the entry's wording, since a different reading of the condition gives a
different digraph and different counts.

Second, the recurrence \eqref{eq:conj} was evaluated on those published terms in exact
integer arithmetic, with no matrices involved, and vanishes wherever it applies.

Third, the annihilation test of Lemma~\ref{lem:crit} was run in exact integer arithmetic
until the working vector was identically zero, not to a sampled prefix, and returned zero
throughout.

\begin{thebibliography}{9}
\bibitem{oeis} The OEIS Foundation, \emph{The On-Line Encyclopedia of Integer
Sequences}, \url{https://oeis.org}, 2026. Sequence A234121.
\bibitem{stanley} R.~P.~Stanley, \emph{Enumerative Combinatorics}, Volume 1, 2nd ed.,
Cambridge University Press, 2011. (Transfer-matrix method, Section 4.7.)
\bibitem{fs} P.~Flajolet and R.~Sedgewick, \emph{Analytic Combinatorics}, Cambridge
University Press, 2009.
\bibitem{horn} R.~A.~Horn and C.~R.~Johnson, \emph{Matrix Analysis}, 2nd ed., Cambridge
University Press, 2013.
\end{thebibliography}

\end{document}
